% GENERATED FILE. Edit canonical lessons, then run npm run generate:books.
% canonical-source-hash: fnv1a64:4a664ff83191bdbf
% canonical-lessons: FR-C444-charcuterie, FR-C444-poissonnerie, FR-C444-traiteur, FR-C444-epicier, FR-C444-primeur

\chapter{Les commerces de bouche}
\label{ch:fr-les-commerces-de-bouche}
\hlchaptermodality{444}

\begin{chapteropening}
\textbf{By the end of this chapter:} \emph{I can shop for food in small shops: the delicatessen, the fishmonger, the caterer, the grocer and the greengrocer.}

Complete the last lesson of chapter 444: I can shop for food in small shops: the delicatessen, the fishmonger, the caterer, the grocer and the greengrocer.
\end{chapteropening}

\section[la charcuterie]{la charcuterie — the pork butcher, the delicatessen; cold meats}

\label{lesson:FR-C444-charcuterie}



\begin{quote}
Before the new one: say the French for to salt, then the French for to clear of snow.
\end{quote}

\subsection*{You'll want to know: la charcuterie}
\textbf{la charcuterie} — \textquotedblleft{}the pork butcher, the delicatessen; cold meats\textquotedblright{}.

\begin{quote}
À la charcuterie, nous achetons du jambon et du pâté.
\end{quote}

\begin{quote}
\emph{At the pork butcher's we buy ham and pâté.}
\end{quote}

\emph{De la charcuterie} also means the cold meats themselves.

\begin{cousinweb}
\emph{charcuterie} is built on \emph{charcutier}, from \emph{chair cuite} 'cooked meat'; \emph{chair} goes back to Latin \emph{carō} 'flesh' and \emph{cuite} to \emph{coquere} 'to cook' (Lewis \& Short). EtymDB lists it without a source.
\end{cousinweb}

\begin{grammarlens}[title={What you've built}]
One more word for telling, explaining and describing at B1.
\end{grammarlens}

\subsection*{Your turn}
\begin{itemize}
\raggedright
  \item \emph{Say it:} \emph{la charcuterie}
  \item \emph{Say it:} \emph{la charcuterie}, once more
  \item \emph{Say it:} a sentence of your own with \emph{la charcuterie}
  \item \emph{Recall:} say the French for to salt, then the French for to clear of snow, then say \emph{la charcuterie} again
\end{itemize}

\begin{tcolorbox}[breakable,colback=teal!4,colframe=teal!35!black,title={Before you move on}]
What does la charcuterie mean? (\textquotedblleft{}The pork butcher, the delicatessen; cold meats\textquotedblright{}.) Say it once more.
\end{tcolorbox}
\section[la poissonnerie]{la poissonnerie — the fishmonger’s, the fish shop}

\label{lesson:FR-C444-poissonnerie}



\begin{quote}
Before the new one: say the French for to clear of snow, then the French for the pork butcher.
\end{quote}

\subsection*{You'll want to know: la poissonnerie}
\textbf{la poissonnerie} — \textquotedblleft{}the fishmonger’s, the fish shop\textquotedblright{}.

\begin{quote}
Le vendredi, il y a la queue devant la poissonnerie.
\end{quote}

\begin{quote}
\emph{On Fridays there is a queue outside the fish shop.}
\end{quote}

\emph{Le poissonnier} — the fishmonger.

\begin{cousinweb}
\emph{poissonnerie} is built on \emph{poisson} 'fish', which goes back to Latin \emph{piscis} (Lewis \& Short). EtymDB does not list the shop name.
\end{cousinweb}

\begin{grammarlens}[title={What you've built}]
One more word for telling, explaining and describing at B1.
\end{grammarlens}

\subsection*{Your turn}
\begin{itemize}
\raggedright
  \item \emph{Say it:} \emph{la poissonnerie}
  \item \emph{Say it:} \emph{la poissonnerie}, once more
  \item \emph{Say it:} a sentence of your own with \emph{la poissonnerie}
  \item \emph{Recall:} say the French for to clear of snow, then the French for the pork butcher, then say \emph{la poissonnerie} again
\end{itemize}

\begin{tcolorbox}[breakable,colback=teal!4,colframe=teal!35!black,title={Before you move on}]
What does la poissonnerie mean? (\textquotedblleft{}The fishmonger’s, the fish shop\textquotedblright{}.) Say it once more.
\end{tcolorbox}
\section[le traiteur]{le traiteur — the caterer; the deli (with ready meals)}

\label{lesson:FR-C444-traiteur}



\begin{quote}
Before the new one: say the French for the pork butcher, then the French for the fishmonger’s.
\end{quote}

\subsection*{You'll want to know: le traiteur}
\textbf{le traiteur} — \textquotedblleft{}the caterer; the deli (with ready meals)\textquotedblright{}.

\begin{quote}
Pour la fête, nous avons commandé des plats chez le traiteur.
\end{quote}

\begin{quote}
\emph{For the party we ordered dishes from the caterer.}
\end{quote}

A shop that sells ready-made dishes and caters for parties.

\begin{cousinweb}
\emph{traiteur} is built on \emph{traiter} 'to deal with, to serve food to', which is inherited from Latin \emph{tractāre} 'to handle', from \emph{trahere} 'to draw' (Lewis \& Short).
\end{cousinweb}

\begin{grammarlens}[title={What you've built}]
One more word for telling, explaining and describing at B1.
\end{grammarlens}

\subsection*{Your turn}
\begin{itemize}
\raggedright
  \item \emph{Say it:} \emph{le traiteur}
  \item \emph{Say it:} \emph{le traiteur}, once more
  \item \emph{Say it:} a sentence of your own with \emph{le traiteur}
  \item \emph{Recall:} say the French for the pork butcher, then the French for the fishmonger’s, then say \emph{le traiteur} again
\end{itemize}

\begin{tcolorbox}[breakable,colback=teal!4,colframe=teal!35!black,title={Before you move on}]
What does le traiteur mean? (\textquotedblleft{}The caterer; the deli (with ready meals)\textquotedblright{}.) Say it once more.
\end{tcolorbox}
\section[l'épicier]{l'épicier — the grocer}

\label{lesson:FR-C444-epicier}



\begin{quote}
Before the new one: say the French for the fishmonger’s, then the French for the caterer.
\end{quote}

\subsection*{You'll want to know: l'épicier}
\textbf{l'épicier} — \textquotedblleft{}the grocer\textquotedblright{}.

\begin{quote}
L'épicier du quartier ferme à neuf heures.
\end{quote}

\begin{quote}
\emph{The local grocer closes at nine.}
\end{quote}

\emph{L'épicerie} — the grocer's shop.

\begin{cousinweb}
\emph{épicier} is built on \emph{épice} 'spice', which is inherited from Latin \emph{speciēs} 'kind', later 'goods, spices' (Lewis \& Short).
\end{cousinweb}

\begin{grammarlens}[title={What you've built}]
One more word for telling, explaining and describing at B1.
\end{grammarlens}

\subsection*{Your turn}
\begin{itemize}
\raggedright
  \item \emph{Say it:} \emph{l'épicier}
  \item \emph{Say it:} \emph{l'épicier}, once more
  \item \emph{Say it:} a sentence of your own with \emph{l'épicier}
  \item \emph{Recall:} say the French for the fishmonger’s, then the French for the caterer, then say \emph{l'épicier} again
\end{itemize}

\begin{tcolorbox}[breakable,colback=teal!4,colframe=teal!35!black,title={Before you move on}]
What does l'épicier mean? (\textquotedblleft{}The grocer\textquotedblright{}.) Say it once more.
\end{tcolorbox}
\section[le primeur]{le primeur — the greengrocer; early produce}

\label{lesson:FR-C444-primeur}



\begin{quote}
Before the new one: say the French for the caterer, then the French for the grocer.
\end{quote}

\subsection*{You'll want to know: le primeur}
\textbf{le primeur} — \textquotedblleft{}the greengrocer; early produce\textquotedblright{}.

\begin{quote}
Chez le primeur, les fraises sont déjà arrivées.
\end{quote}

\begin{quote}
\emph{At the greengrocer's the strawberries are already in.}
\end{quote}

\emph{Les primeurs} are the first fruit and vegetables of the year.

\begin{cousinweb}
\emph{primeur} is built on the old adjective \emph{prime} 'first', from Latin \emph{prīmus} (Lewis \& Short): the first produce of the year. EtymDB does not list it.
\end{cousinweb}

\begin{grammarlens}[title={What you've built}]
One more word for telling, explaining and describing at B1.
\end{grammarlens}

\subsection*{Your turn}
\begin{itemize}
\raggedright
  \item \emph{Say it:} \emph{le primeur}
  \item \emph{Say it:} \emph{le primeur}, once more
  \item \emph{Say it:} a sentence of your own with \emph{le primeur}
  \item \emph{Recall:} say the French for the caterer, then the French for the grocer, then say \emph{le primeur} again
\end{itemize}

\begin{tcolorbox}[breakable,colback=teal!4,colframe=teal!35!black,title={Before you move on}]
What does le primeur mean? (\textquotedblleft{}The greengrocer; early produce\textquotedblright{}.) Say it once more.
\end{tcolorbox}
